\documentclass[11pt,a4paper]{article}
\usepackage[margin=1in]{geometry}
\usepackage{amsmath,amssymb,amsthm}
\usepackage{algorithm}
\usepackage{algpseudocode}
\usepackage{booktabs}
\usepackage{hyperref}

\theoremstyle{definition}
\newtheorem{definition}{Definition}[section]
\newtheorem{theorem}{Theorem}[section]
\newtheorem{lemma}[theorem]{Lemma}
\newtheorem{remark}{Remark}[section]

\title{\textbf{Mathematical Foundations, Attention Logit Bounding, and Gradient Saturation Dynamics of QK-Norm and Logit Soft-Capping}}
\author{Team Scaibu \\ \small Email: \texttt{jaydeep.9664920749@gmail.com} \\ \small Architecture and Core Engine Team}
\date{October 2026}

\begin{document}
\maketitle

\begin{abstract}
This document provides the rigorous mathematical formulation, variance bounding, and numerical stability analysis for Query-Key Normalization (QK-Norm) and Logit Soft-Capping in modern large language models. Unconstrained dot-product attention logits frequently grow with sequence depth, leading to attention entropy collapse and loss spikes during large-scale training. QK-Norm bounds the pre-attention vector norms on the unit hypersphere, while hyperbolic tangent soft-capping smoothly compresses logits into $(-c, c)$. We prove strict boundedness properties, derive operational complexities, trace a complete worked numerical example, and document kernel integration considerations.
\end{abstract}

\section{Notation and Mathematical Preliminaries}

Let $Q \in \mathbb{R}^{N_q \times D}$ and $K \in \mathbb{R}^{N_k \times D}$ denote query and key activation matrices for an attention head with sequence lengths $N_q, N_k$ and head dimension $D$.

\begin{itemize}
    \item $N_q, N_k \in \mathbb{N}_{\ge 1}$: Query and Key sequence token lengths.
    \item $D \in \mathbb{N}_{\ge 1}$: Head projection dimension ($d_k$).
    \item $q_i \in \mathbb{R}^D$: Query vector for token $i \in \{1, \dots, N_q\}$.
    \item $k_j \in \mathbb{R}^D$: Key vector for token $j \in \{1, \dots, N_k\}$.
    \item $\hat{q}_i, \hat{k}_j \in \mathbb{R}^D$: Normalized query and key representations.
    \item $\tau = \frac{1}{\sqrt{D}}$: Canonical attention scaling temperature factor.
    \item $c \in \mathbb{R}_{> 0}$: Soft-capping threshold asymptote ($c = 50.0$ for attention logits, $c = 30.0$ for vocabulary logits).
    \item $A \in \mathbb{R}^{N_q \times N_k}$: Final capped attention logit matrix.
    \item $\epsilon > 0$: Denominator stabilization constant ($\epsilon = 10^{-6}$).
\end{itemize}

\begin{table}[h]
\centering
\caption{Summary of Mathematical Symbols for QK-Norm and Soft-Capping}
\begin{tabular}{lll}
\toprule
\textbf{Symbol} & \textbf{Type / Domain} & \textbf{Description} \\
\midrule
$Q, K$ & $\mathbb{R}^{N_q \times D}, \mathbb{R}^{N_k \times D}$ & Raw query and key matrices \\
$\hat{Q}, \hat{K}$ & $\mathbb{R}^{N_q \times D}, \mathbb{R}^{N_k \times D}$ & Hypersphere-normalized queries and keys \\
$S$ & $\mathbb{R}^{N_q \times N_k}$ & Raw scaled dot-product matrix \\
$A$ & $(-c, c)^{N_q \times N_k}$ & Soft-capped attention logit tensor \\
$c$ & $\mathbb{R}_{> 0}$ & Soft-capping bound parameter ($50.0$) \\
$\tau$ & $\mathbb{R}_{> 0}$ & Attention scaling factor ($1/\sqrt{D}$) \\
$\epsilon$ & $\mathbb{R}_{> 0}$ & Regularization constant ($10^{-6}$) \\
\bottomrule
\end{tabular}
\end{table}

\section{Problem Definition and Core Concepts}

\subsection{Intuitive Meaning}
In standard attention $A_{ij} = \frac{q_i^\top k_j}{\sqrt{D}}$, as models grow deeper and learning rates increase, the magnitudes $\|q_i\|$ and $\|k_j\|$ can diverge, causing $A_{ij}$ to reach extreme values ($> 100$). When passed through Softmax, large logits cause gradient vanishing (one-hot collapse) or numerical overflow in FP16. QK-Norm projects vectors onto a fixed-norm sphere, and soft-capping smoothly restricts logits using a hyperbolic tangent envelope.

\subsection{Formal Mathematical Formulation}
\begin{align}
\hat{q}_i &= \text{RMSNorm}(q_i) = \frac{q_i}{\sqrt{\frac{1}{D} \sum_{d=1}^D q_{id}^2 + \epsilon}} \odot \gamma \\
\hat{k}_j &= \text{RMSNorm}(k_j) = \frac{k_j}{\sqrt{\frac{1}{D} \sum_{d=1}^D k_{jd}^2 + \epsilon}} \odot \gamma \\
S_{ij} &= \tau \cdot (\hat{q}_i^\top \hat{k}_j) = \frac{1}{\sqrt{D}} \sum_{d=1}^D \hat{q}_{id} \hat{k}_{jd} \\
A_{ij} &= c \cdot \tanh\left( \frac{S_{ij}}{c} \right)
\end{align}

\section{Algorithm Specification}

\begin{algorithm}[H]
\caption{QK-Norm and Attention Logit Soft-Capping}
\label{alg:qk_norm_soft_cap}
\begin{algorithmic}[1]
\Require Queries $Q \in \mathbb{R}^{N_q \times D}$, Keys $K \in \mathbb{R}^{N_k \times D}$, scale $\tau > 0$, cap threshold $c > 0$, regularizer $\epsilon > 0$.
\Ensure Logits $A \in \mathbb{R}^{N_q \times N_k}$, normalized $\hat{Q} \in \mathbb{R}^{N_q \times D}$, $\hat{K} \in \mathbb{R}^{N_k \times D}$.
\For{$i \gets 1 \textbf{ to } N_q$}
    \State $\text{rms}_q \gets \sqrt{\frac{1}{D}\sum_{d=1}^D Q_{id}^2 + \epsilon}$
    \State $\hat{Q}_{i, :} \gets Q_{i, :} / \text{rms}_q$
\EndFor
\For{$j \gets 1 \textbf{ to } N_k$}
    \State $\text{rms}_k \gets \sqrt{\frac{1}{D}\sum_{d=1}^D K_{jd}^2 + \epsilon}$
    \State $\hat{K}_{j, :} \gets K_{j, :} / \text{rms}_k$
\EndFor
\For{$i \gets 1 \textbf{ to } N_q$}
    \For{$j \gets 1 \textbf{ to } N_k$}
        \State $S_{ij} \gets \tau \sum_{d=1}^D \hat{Q}_{id} \hat{K}_{jd}$
        \State $A_{ij} \gets c \cdot \tanh(S_{ij} / c)$
    \EndFor
\EndFor
\State \Return $\{A, \hat{Q}, \hat{K}\}$
\end{algorithmic}
\end{algorithm}

\section{Theoretical Guarantees and Logit Bounds}

\begin{theorem}[Strict Upper Bounding of Attention Logits]
For any unconstrained query $q \in \mathbb{R}^D$ and key $k \in \mathbb{R}^D$, under QK-Norm with unit gain $\gamma = \mathbf{1}$ and soft-capping threshold $c > 0$:
\begin{enumerate}
    \item $\|\hat{q}\|_2 = \sqrt{D}$ and $\|\hat{k}\|_2 = \sqrt{D}$.
    \item The un-capped scaled dot product satisfies $|S_{ij}| \le \sqrt{D}$.
    \item The soft-capped logit strictly satisfies:
    \begin{equation}
    -c < A_{ij} < c
    \end{equation}
\end{enumerate}
\end{theorem}

\begin{proof}
By definition of RMS normalization:
\begin{equation}
\|\hat{q}\|_2^2 = \sum_{d=1}^D \frac{q_d^2}{\frac{1}{D} \sum_{m=1}^D q_m^2} = D \frac{\sum q_d^2}{\sum q_m^2} = D \implies \|\hat{q}\|_2 = \sqrt{D}
\end{equation}
By the Cauchy-Schwarz inequality:
\begin{equation}
|S_{ij}| = \left| \frac{1}{\sqrt{D}} \hat{q}^\top \hat{k} \right| \le \frac{1}{\sqrt{D}} \|\hat{q}\|_2 \|\hat{k}\|_2 = \frac{1}{\sqrt{D}} (\sqrt{D})(\sqrt{D}) = \sqrt{D}
\end{equation}
Finally, because the range of $\tanh(z)$ is the open interval $(-1, 1)$ for all $z \in \mathbb{R}$, $A_{ij} = c \tanh(S_{ij}/c)$ is strictly bounded within $(-c, c)$.
\end{proof}

\subsection{Limitations}
\begin{itemize}
    \item \textbf{Kernel Non-Linearity}: Hyperbolic tangent soft-capping prevents standard fused FlashAttention-2 kernels from using pure matrix accumulators without specialized kernel extensions.
\end{itemize}

\section{Complexity Analysis}

\subsection{Time Complexity}
\begin{itemize}
    \item \textbf{Normalization}: Normalizing $Q$ takes $2 N_q D$ FLOPs; normalizing $K$ takes $2 N_k D$ FLOPs.
    \item \textbf{Dot-Product and Capping}: Computing $S$ takes $2 N_q N_k D$ FLOPs; tanh mapping takes $\mathcal{O}(N_q N_k)$ operations.
    \item \textbf{Total Time Complexity}: $\Theta(N_q N_k D)$ FLOPs.
\end{itemize}

\subsection{Space Complexity}
\begin{itemize}
    \item \textbf{Activation Storage}: Requires storing normalized $\hat{Q}, \hat{K}$ ($\mathcal{O}((N_q + N_k)D)$) and logits $A$ ($\mathcal{O}(N_q N_k)$).
\end{itemize}

\section{Exactness and Error Analysis}

The operation is mathematically exact up to floating-point precision. For small input arguments ($|S_{ij}/c| \ll 1$), the Taylor expansion of $\tanh(z) = z - \frac{z^3}{3} + \dots$ yields $c \tanh(S/c) \approx S$, ensuring that soft-capping preserves linear attention dynamics in normal operating ranges while strictly preventing extreme outlier divergence.

\section{Worked Numerical Example}

Let $N_q = 1, N_k = 1, D = 2, c = 10.0, \epsilon = 10^{-6}$.
\begin{equation}
Q = [[3.0, 4.0]], \quad K = [[6.0, 8.0]]
\end{equation}

\begin{enumerate}
    \item \textbf{Query RMS Normalization}:
    \begin{align}
    \text{MS}_q &= \frac{3^2 + 4^2}{2} = \frac{9 + 16}{2} = 12.5 \\
    \text{RMS}_q &= \sqrt{12.5 + 10^{-6}} \approx 3.535534 \\
    \hat{Q} &= \left[ \frac{3.0}{3.535534}, \frac{4.0}{3.535534} \right] \approx [0.848528, 1.131371]
    \end{align}
    \item \textbf{Key RMS Normalization}:
    \begin{align}
    \text{MS}_k &= \frac{6^2 + 8^2}{2} = \frac{36 + 64}{2} = 50.0 \\
    \text{RMS}_k &= \sqrt{50.0 + 10^{-6}} \approx 7.071068 \\
    \hat{K} &= \left[ \frac{6.0}{7.071068}, \frac{8.0}{7.071068} \right] \approx [0.848528, 1.131371]
    \end{align}
    \item \textbf{Dot Product and Temperature Scaling} ($\tau = 1/\sqrt{2} \approx 0.7071068$):
    \begin{align}
    \text{Dot} &= 0.848528^2 + 1.131371^2 = 0.720000 + 1.280000 = 2.000000 \\
    S &= 2.0 \times 0.7071068 = 1.4142136
    \end{align}
    \item \textbf{Soft-Capped Logit} ($c = 10.0$):
    \begin{equation}
    A = 10.0 \times \tanh\left( \frac{1.4142136}{10.0} \right) = 10.0 \times \tanh(0.14142136) \approx 10.0 \times 0.140486 = 1.404860
    \end{equation}
\end{enumerate}

\section{Numerical Considerations and Edge Cases}

\begin{itemize}
    \item \textbf{Asymptotic Gradient}: The derivative $\frac{d}{dS}[c \tanh(S/c)] = 1 - \tanh^2(S/c) = \text{sech}^2(S/c)$. As $S \to \infty$, gradients smoothly diminish to zero, automatically suppressing disruptive gradient spikes from outlier tokens.
\end{itemize}

\begin{thebibliography}{9}
\bibitem{deghani2023scaling} M.~Dehghani et al., ``Scaling Vision Transformers to 22 Billion Parameters,'' in \emph{International Conference on Machine Learning (ICML)}, pp.~7480--7498, 2023.
\bibitem{gemma2024technical} Gemma Team, Google DeepMind, ``Gemma 2: Improving Open Language Models at a Practical Size,'' \emph{arXiv preprint arXiv:2408.00118}, 2024.
\end{thebibliography}

\end{document}
